% \iffalse meta-comment
%
% Copyright (C) SJTUG
%   2018--2026 Weijian Wu   <alexarawu@outlook.com>
%   2022--2026 Zilong Li    <logcreative@outlook.com>
%   2024--2026 Boshi Yuan   <nemoyuan2008@outlook.com>
%
% This work may be distributed and/or modified under the
% conditions of the LaTeX Project Public License, either version 1.3c
% of this license or (at your option) any later version.
% The latest version of this license is in
%   https://www.latex-project.org/lppl.txt
% and version 1.3c or later is part of all distributions of LaTeX
% version 2008 or later.
%
%<*!(driver|install)>
%<+!driver>\GetIdInfo $Id: sjtutex-lang-scheme.dtx 2.3.1 2026-04-21 15:09:02 +0800 f6eb9dc Alexara Wu <alexarawu@outlook.com>$
%<lang&zh>  {Chinese specific definition (SJTUTeX)}
%<lang&zh>\ProvidesExplFile{sjtu-lang-zh.def}
%<lang&en>  {English specific definition (SJTUTeX)}
%<lang&en>\ProvidesExplFile{sjtu-lang-en.def}
%<lang&de>  {German specific definition (SJTUTeX)}
%<lang&de>\ProvidesExplFile{sjtu-lang-de.def}
%<lang&ja>  {Japanese specific definition (SJTUTeX)}
%<lang&ja>\ProvidesExplFile{sjtu-lang-ja.def}
%<scheme&zh>  {Chinese scheme (SJTUTeX)}
%<scheme&zh>\ProvidesExplFile{sjtu-scheme-zh.def}
%<scheme&en>  {English scheme (SJTUTeX)}
%<scheme&en>\ProvidesExplFile{sjtu-scheme-en.def}
%<scheme&de>  {German scheme (SJTUTeX)}
%<scheme&de>\ProvidesExplFile{sjtu-scheme-de.def}
%<scheme&ja>  {Japanese scheme (SJTUTeX)}
%<scheme&ja>\ProvidesExplFile{sjtu-scheme-ja.def}
%<!driver>  {\ExplFileDate}{\ExplFileVersion}{\ExplFileDescription}
%</!(driver|install)>
%
% \fi
%
% \begin{implementation}
%
% \paragraph*{\fbox{\file{sjtutex-lang-scheme.dtx}}}
%
% 本文件定义各语种的名称、日期与版式设定（语言支持与方案）。
% 语言与方案配置文件在文档类载入时按语言读取，
% 其实现见 \file{sjtutex.dtx} 中的多语言支持小节。
%
% \subsection{语言与方案}
%
%    \begin{macrocode}
%<*scheme>
%<@@=sjtu>
%<zh|ja>\keys_set_known:nn { ctex / chapter }
%<zh|ja>  { name = { 第 \space , \space 章 } }
%</scheme>
%    \end{macrocode}
%
% 设置全文首行缩进。
%    \begin{macrocode}
%<*scheme>
\ctex_if_autoindent_touched:F
%<zh>  { \ctex_set:n { autoindent = true } }
%<en|de>  { \ctex_set:n { autoindent = 1.5 em } }
%<ja>  { \ctex_set:n { autoindent = 1 } }
%    \end{macrocode}
%
% \begin{macro}{\verse,\quotation}
% 修改诗歌和引用环境的缩进。
%    \begin{macrocode}
%<zh|ja>\ctex_patch_cmd:Nnn \verse { -1.5em } { -2 \ccwd }
%<zh|ja>\ctex_patch_cmd:Nnn \verse {  1.5em } {  2 \ccwd }
\ctex_patch_cmd:Nnn \quotation { 1.5em } { \parindent }
%</scheme>
%    \end{macrocode}
% \end{macro}
%
% 设置脚注标记样式。
%    \begin{macrocode}
%<*scheme>
\keys_set:nn { sjtu / style }
%<zh|ja>  { fnmark-style = circled }
%<en|de>  { fnmark-style = plain   }
%</scheme>
%    \end{macrocode}
%
% \begin{macro}[int]{\@@_set_cjk_default_zh:,\@@_set_cjk_default_ja:}
% 设置 CJK 默认字体族的辅助命令。
%    \begin{macrocode}
%<*lang>
%<*zh>
\cs_new_protected:Nn \@@_set_cjk_default_zh:
  {
    \tl_set:Nn \CJKrmdefault { zhsong }
    \tl_set:Nn \CJKsfdefault { zhhei  }
    \tl_set:Nn \CJKttdefault { zhfs   }
  }
%</zh>
%<*ja>
\cs_new_protected:Nn \@@_set_cjk_default_ja:
  {
    \tl_set:Nn \CJKrmdefault { jamin  }
    \tl_set:Nn \CJKsfdefault { jagoth }
    \tl_set:Nn \CJKttdefault { jagoth }
  }
%</ja>
%    \end{macrocode}
% \end{macro}
%
% \changes{v2.2.1}{2025/03/27}{新增语言设置钩子。}
% 语言设置钩子。
%    \begin{macrocode}
%<zh>\NewHook { sjtutex / lang / zh }
%<en>\NewHook { sjtutex / lang / en }
%<de>\NewHook { sjtutex / lang / de }
%<ja>\NewHook { sjtutex / lang / ja }
%    \end{macrocode}
%
% 定义语言切换键，设置 CJK 默认字体族并触发语言钩子。
%    \begin{macrocode}
\keys_define:nn { sjtu / private }
  {
%<zh>    lang / zh .code:n =
%<en>    lang / en .code:n =
%<de>    lang / de .code:n =
%<ja>    lang / ja .code:n =
      {
        \tl_set_eq:NN \l_@@_lang_tl \l_keys_value_tl
%<zh>        \@@_set_cjk_default_zh:
%<ja>        \@@_set_cjk_default_ja:
        \normalfont
%<zh>        \UseHook { sjtutex / lang / zh }
%<en>        \UseHook { sjtutex / lang / en }
%<de>        \UseHook { sjtutex / lang / de }
%<ja>        \UseHook { sjtutex / lang / ja }
      } ,
%<zh>    zh .meta:n = { lang = zh } ,
%<zh>    zh .groups:n = { lang }
%<en>    en .meta:n = { lang = en } ,
%<en>    en .groups:n = { lang }
%<de>    de .meta:n = { lang = de } ,
%<de>    de .groups:n = { lang }
%<ja>    ja .meta:n = { lang = ja } ,
%<ja>    ja .groups:n = { lang }
  }
%    \end{macrocode}
%
% 各语言钩子的默认设置。
%    \begin{macrocode}
%<zh>\AddToHook { sjtutex / lang / zh }
%<en>\AddToHook { sjtutex / lang / en }
%<de>\AddToHook { sjtutex / lang / de }
%<ja>\AddToHook { sjtutex / lang / ja }
  {
%<zh>    \tl_set:Nn \languagename { chinese }
%<en>    \tl_set:Nn \languagename { english }
%<de>    \tl_set:Nn \languagename { ngerman }
%<ja>    \tl_set:Nn \languagename { japanese }
%<zh>    \ctex_set:n { autoindent = true }
%<en|de>    \ctex_set:n { autoindent = 1.5 em }
%<ja>    \ctex_set:n { autoindent = 1 }
  }
%    \end{macrocode}
%
% \begin{macro}[int]{\@@_title_case_aux_zh:n,\@@_title_case_aux_en:n,
% \@@_title_case_aux_de:n,\@@_title_case_aux_ja:n}
% 设置标题大小写转换的辅助命令。
%    \begin{macrocode}
%<zh>\cs_set_eq:NN \@@_title_case_aux_zh:n \use:n
%<en>\cs_set_eq:NN \@@_title_case_aux_en:n \MakeUppercase
%<de>\cs_set_eq:NN \@@_title_case_aux_de:n \MakeUppercase
%<ja>\cs_set_eq:NN \@@_title_case_aux_ja:n \use:n
%    \end{macrocode}
% \end{macro}
%
% 将形如 |yyyy-mm-dd| 或 |yyyy-mm| 的 ISO 日期格式字符串转化为日期表示。
%
% 日期常量。
%    \begin{macrocode}
%<zh>\@@_name_const:nn { zh } { year = 年, month = 月, day = 日 }
%<ja>\@@_name_const:nn { ja } { year = 年, month = 月, day = 日 }
%<*en>
\clist_const:Nn \c_@@_name_month_en_clist
  {
    January, February, March, April, May, June,
    July, August, September, October, November, December
  }
%</en>
%<*de>
\clist_const:Nn \c_@@_name_month_de_clist
  {
    Januar, Februar, März, April, Mai, Juni,
    Juli, August, September, Oktober, November, Dezember
  }
%</de>
%    \end{macrocode}
%
% \begin{macro}[int]{\@@_date_aux_zh:nnn,\@@_date_aux_zh:w,
% \@@_date_aux_short_zh:nn,\@@_date_aux_short_zh:w}
% 中文日期。
%    \begin{macrocode}
%<*zh>
\cs_new:Npn \@@_date_aux_zh:nnn #1#2#3
  {
    \int_to_arabic:n {#1} ~ { \exp_not:V \c_@@_name_year_zh_tl  } ~
    \int_to_arabic:n {#2} ~ { \exp_not:V \c_@@_name_month_zh_tl } ~
    \int_to_arabic:n {#3} ~ { \exp_not:V \c_@@_name_day_zh_tl   }
  }
\cs_new:Npn \@@_date_aux_zh:w #1-#2-#3 \q_stop
  { \@@_date_aux_zh:nnn {#1} {#2} {#3} }
\cs_new:Npn \@@_date_aux_short_zh:nn #1#2
  {
    \int_to_arabic:n {#1} ~ { \exp_not:V \c_@@_name_year_zh_tl  } ~
    \int_to_arabic:n {#2} ~ { \exp_not:V \c_@@_name_month_zh_tl }
  }
\cs_new:Npn \@@_date_aux_short_zh:w #1-#2 \q_stop
  { \@@_date_aux_short_zh:nn {#1} {#2} }
%</zh>
%    \end{macrocode}
% \end{macro}
%
% \begin{macro}[int]{\@@_ordinal_en:n}
% 上标形式的序数词。
%    \begin{macrocode}
%<*en>
\cs_new:Npn \@@_ordinal_en:n #1
  {
    \int_to_arabic:n {#1}
    \exp_not:N \textsuperscript
      {
        \int_case:nnF { \int_mod:nn {#1} { 100 } }
          {
            { 11 } { th }
            { 12 } { th }
            { 13 } { th }
          }
          {
            \int_case:nnF { \int_mod:nn {#1} { 10 } }
              {
                { 1 } { st }
                { 2 } { nd }
                { 3 } { rd }
              }
              { th }
          }
      }
  }
%    \end{macrocode}
% \end{macro}
%
% \begin{macro}[int]{\@@_date_aux_en:nnn,\@@_date_aux_en:w,
% \@@_date_aux_short_en:nn,\@@_date_aux_short_en:w}
% 英文日期。
%    \begin{macrocode}
\cs_new:Npn \@@_date_aux_en:nnn #1#2#3
  {
    \clist_item:Nn \c_@@_name_month_en_clist {#2} ~
    \@@_ordinal_en:n {#3} ,~
    \int_to_arabic:n {#1}
  }
\cs_new:Npn \@@_date_aux_en:w #1-#2-#3 \q_stop
  { \@@_date_aux_en:nnn {#1} {#2} {#3} }
\cs_new:Npn \@@_date_aux_short_en:nn #1#2
  {
    \clist_item:Nn \c_@@_name_month_en_clist {#2} ,~
    \int_to_arabic:n {#1}
  }
\cs_new:Npn \@@_date_aux_short_en:w #1-#2 \q_stop
  { \@@_date_aux_short_en:nn {#1} {#2} }
%</en>
%    \end{macrocode}
% \end{macro}
%
% \begin{macro}[int]{\@@_date_aux_de:nnn,\@@_date_aux_de:w,
% \@@_date_aux_short_de:nn,\@@_date_aux_short_de:w}
% 德文日期。
%    \begin{macrocode}
%<*de>
\cs_new:Npn \@@_date_aux_de:nnn #1#2#3
  {
    \clist_item:Nn \c_@@_name_month_de_clist {#2} ~
    {#3} ,~ \int_to_arabic:n {#1}
  }
\cs_new:Npn \@@_date_aux_de:w #1-#2-#3 \q_stop
  { \@@_date_aux_de:nnn {#1} {#2} {#3} }
\cs_new:Npn \@@_date_aux_short_de:nn #1#2
  {
    \clist_item:Nn \c_@@_name_month_de_clist {#2} ,~
    \int_to_arabic:n {#1}
  }
\cs_new:Npn \@@_date_aux_short_de:w #1-#2 \q_stop
  { \@@_date_aux_short_de:nn {#1} {#2} }
%</de>
%    \end{macrocode}
% \end{macro}
%
% \begin{macro}[int]{\@@_date_aux_ja:nnn,\@@_date_aux_ja:w,
% \@@_date_aux_short_ja:nn,\@@_date_aux_short_ja:w}
% 日文日期。
%    \begin{macrocode}
%<*ja>
\cs_new:Npn \@@_date_aux_ja:nnn #1#2#3
  {
    \int_to_arabic:n {#1} ~ { \exp_not:V \c_@@_name_year_ja_tl  } ~
    \int_to_arabic:n {#2} ~ { \exp_not:V \c_@@_name_month_ja_tl } ~
    \int_to_arabic:n {#3} ~ { \exp_not:V \c_@@_name_day_ja_tl   }
  }
\cs_new:Npn \@@_date_aux_ja:w #1-#2-#3 \q_stop
  { \@@_date_aux_ja:nnn {#1} {#2} {#3} }
\cs_new:Npn \@@_date_aux_short_ja:nn #1#2
  {
    \int_to_arabic:n {#1} ~ { \exp_not:V \c_@@_name_year_ja_tl  } ~
    \int_to_arabic:n {#2} ~ { \exp_not:V \c_@@_name_month_ja_tl }
  }
\cs_new:Npn \@@_date_aux_short_ja:w #1-#2 \q_stop
  { \@@_date_aux_short_ja:nn {#1} {#2} }
%</ja>
%</lang>
%    \end{macrocode}
% \end{macro}
%
% 初始化语言名称。
%    \begin{macrocode}
%<*scheme>
%<zh>\tl_set:Nn \languagename { chinese  }
%<en>\tl_set:Nn \languagename { english  }
%<de>\tl_set:Nn \languagename { ngerman  }
%<ja>\tl_set:Nn \languagename { japanese }
%    \end{macrocode}
%
% 设置名称选项。默认值为英文，只需修改其他语种。
%    \begin{macrocode}
%<*!en>
\keys_set_known:nn { sjtu / name }
  {
%<*zh>
    contents      = { 目 \protect \quad 录   } ,
    listfigure    = { 插 \protect \quad 图   } ,
    listtable     = { 表 \protect \quad 格   } ,
    figure        = { 图                     } ,
    table         = { 表                     } ,
    abstract      = { 摘 \protect \quad 要   } ,
    index         = { 索 \protect \quad 引   } ,
    appendix      = { 附录                   } ,
    proof         = { 证明                   } ,
    bib           = { 参考文献               } ,
    figure*       = { Figure                 } ,
    table*        = { Table                  } ,
    algorithm     = { 算法                   } ,
    listalgorithm = { 算 \protect \quad 法   } ,
    abbr          = { 缩略语对照表           } ,
    nom           = { 符号对照表             } ,
    ack           = { 致 \protect \quad 谢   } ,
    resume        = { 个人简历               } ,
    digest        = { 大摘要                 } ,
    achv          = { 学术论文和科研成果目录 }
%</zh>
%<*de>
    contents      = { Inhaltsverzeichnis     } ,
    listfigure    = { Abbildungsverzeichnis  } ,
    listtable     = { Tabellenverzeichnis    } ,
    figure        = { Abbildung              } ,
    table         = { Tabelle                } ,
    abstract      = { Zusammenfassung        } ,
    index         = { Index                  } ,
    appendix      = { Anhang                 } ,
    proof         = { Beweis                 } ,
    bib           = { Literaturverzeichnis   } ,
    part          = { Teil                   } ,
    chapter       = { Kapitel                } ,
    figure*       = { Figure                 } ,
    table*        = { Table                  } ,
    algorithm     = { Algorithmus            } ,
    listalgorithm = { Algorithmenverzeichnis } ,
    abbr          = { Abkürzungsverzeichnis  } ,
    nom           = { Symbolverzeichnis      } ,
    ack           = { Danksagungen           } ,
    resume        = { Lebenslauf             } ,
    digest        = { Kurzfassung            } ,
    achv          = { Forschungsleistungen   }
%</de>
%<*ja>
    contents      = { 目 \protect \quad 次 } ,
    listfigure    = { 図目次               } ,
    listtable     = { 表目次               } ,
    figure        = { 図                   } ,
    table         = { 表                   } ,
    abstract      = { 概 \protect \quad 要 } ,
    index         = { 索 \protect \quad 引 } ,
    appendix      = { 付録                 } ,
    proof         = { 证明                 } ,
    bib           = { 参考文献             } ,
    figure*       = { Figure               } ,
    table*        = { Table                } ,
    algorithm     = { アルゴリズム         } ,
    listalgorithm = { アルゴリズム目次     } ,
    abbr          = { 略語表               } ,
    nom           = { 記号表               } ,
    ack           = { 謝 \protect \quad 辞 } ,
    resume        = { 履歴書               } ,
    digest        = { 要 \protect \quad 約 } ,
    achv          = { 研究業績書           }
%</ja>
  }
%</!en>
%</scheme>
%    \end{macrocode}
%
% 通用名称常量。
%    \begin{macrocode}
%<*lang>
%<zh>\@@_name_const:nn { zh }
%<en>\@@_name_const:nn { en }
%<de>\@@_name_const:nn { de }
%<ja>\@@_name_const:nn { ja }
  {
%<zh>    keywords = 关键词,
%<en>    keywords = Key~words,
%<de>    keywords = Schlüsselwörter,
%<ja>    keywords = キーワード,
%<zh|ja>    info_sep = ： \null,
%<en|de>    info_sep = \hbox { :~ },
%<zh>    item_sep = ，
%<en|de>    item_sep = { ,~ }
%<ja>    item_sep = \quad
  }
%    \end{macrocode}
%
% 学位论文专用的符号与摘要名称常量，仅 \cls{sjtuthesis} 使用。
%    \begin{macrocode}
\IfClassLoadedF { sjtuthesis } { \endinput }
%<zh>\@@_symbol_const:nn { white_square } { "25A1 }
%<zh>\@@_name_const:nn { zh } { abstract = 摘 \protect \quad 要 }
%<en>\@@_name_const:nn { en } { abstract = Abstract }
%<de>\@@_name_const:nn { de } { abstract = Abstrakt }
%<ja>\@@_name_const:nn { ja } { abstract = 要 \protect \quad 旨 }
%</lang>
%    \end{macrocode}
%
% 预定义的数学环境，不包括证明环境 \env{proof}。
%    \begin{macrocode}
%<*scheme>
%<zh>\@@_name_const:nn { zh }
%<en>\@@_name_const:nn { en }
%<de>\@@_name_const:nn { de }
%<ja>\@@_name_const:nn { ja }
  {
%<*zh>
    assumption  = 假设,
    axiom       = 公理,
    conjecture  = 猜想,
    corollary   = 推论,
    definition  = 定义,
    example     = 例,
    exercise    = 练习,
    lemma       = 引理,
    problem     = 问题,
    proposition = 命题,
    remark      = 注,
    solution    = 解,
    theorem     = 定理
%</zh>
%<*en>
    assumption  = Assumption,
    axiom       = Axiom,
    conjecture  = Conjecture,
    corollary   = Corollary,
    definition  = Definition,
    example     = Example,
    exercise    = Exercise,
    lemma       = Lemma,
    problem     = Problem,
    proposition = Proposition,
    remark      = Remark,
    solution    = Solution,
    theorem     = Theorem
%</en>
%<*de>
    assumption  = Annahme,
    axiom       = Axiom,
    conjecture  = Hypothese,
    corollary   = Korollar,
    definition  = Definition,
    example     = Beispiel,
    exercise    = Übung,
    lemma       = Lemma,
    problem     = Problem,
    proposition = Proposition,
    remark      = Anmerkung,
    solution    = Lösung,
    theorem     = Theorem
%</de>
%<*ja>
    assumption  = 仮定,
    axiom       = 公理,
    conjecture  = 予想,
    corollary   = 系,
    definition  = 定義,
    example     = 例,
    exercise    = 練習,
    lemma       = 補題,
    problem     = 問題,
    proposition = 命題,
    remark      = 注意,
    solution    = 解法,
    theorem     = 定理
%</ja>
  }
%</scheme>
%    \end{macrocode}
% \end{implementation}
%
% \endinput
